\section{An individual Selmer estimate: the vanishing in case Z}
\label{sec:case-z}

Case Z enters \texttt{vwin\_\allowbreak{}of\_\allowbreak{}counts} through \texttt{count\allowbreak{}Z} and
\texttt{prop\allowbreak{}Z}. Here the source proves an exact zero. We derive the
estimate for both residue classes of the twist and then pass to the
character and level used in the window.

\subsection{The strict group and its torsion}

Fix an algebraic closure \(\Omega\) of \(\mathbb Q\), let
\(G=\operatorname{Gal}(\Omega/\mathbb Q)\), and write
\(\kappa:G\to\mathbb Z_\ell^\times\) for the cyclotomic character.
Let \(\mathcal O\) be a discrete valuation ring and a finite domain
over \(\mathbb Z_\ell\) as in \texttt{Coef}. Put
\[
  W=\mathcal O\otimes_{\mathbb Z_\ell}(\mathbb Q_\ell/\mathbb Z_\ell),
     \qquad \alpha=\kappa^{-e}.
\]
The action on \(W\) is scalar multiplication by \(\alpha(g)\).
All cohomology uses continuous cocycles and the discrete topology on
\(W\). Denote by \(T_{\ell,e}\) the submodule of \(H^1(G,W)\)
whose classes restrict to zero on every inertia group and on every
decomposition group above \(\ell\). We prove
\begin{equation}
 \ell\ge5,\quad e\ge2,\quad 2\mid e,\quad \ell-1\mid 2e
           \quad\Longrightarrow\quad T_{\ell,e}=0 .
 \label{eq:rh-z-vanish}
\end{equation}
In the source, \(\operatorname{length}_{\mathcal O}T_{\ell,e}\) is
\(\mathrm{Chr.V}(\ell,e,1)\).

Every element of \(W\) is killed by some power of \(\ell\).
Multiplication by a nonzero element of \(\mathcal O\) on \(W\) is
surjective: \(\mathcal O\) is finite and torsion-free over
\(\mathbb Z_\ell\), so \(W\) identifies with
\(\operatorname{Frac}(\mathcal O)/\mathcal O\).
A continuous cocycle on the compact group \(G\) has finite image in
\(W\). A common power of \(\ell\) therefore kills all its values.
Every cohomology class is consequently killed by some power of \(\ell\).
It suffices to show \(T_{\ell,e}[\ell]=0\), since a nonzero element
of order \(\ell^n\) gives one of order \(\ell\) after multiplication by
\(\ell^{n-1}\). We use the coboundary convention
\(d m(g)=(\alpha(g)-1)m\).

\subsection{The trivial residual character}

Suppose first that \(\ell-1\mid e\). For any
\(u\in\mathbb Z_\ell^\times\), reduction modulo \(\ell\) gives
\(u^{-e}\equiv1\bmod\ell\), because
\(\mathbb F_\ell^\times\) has order \(\ell-1\). Thus the action on
\(W[\ell]\) is trivial. For \(x\in T_{\ell,e}[\ell]\), choose a
representative \(z\). Write \(\ell z=dm\) and take \(m'\in W\)
with \(\ell m'=m\). The representative \(f=z-dm'\) satisfies
\(\ell f=0\) pointwise. Its cocycle identity is now the homomorphism
identity \(f(gh)=f(g)+f(h)\). At a prime other than \(\ell\),
the cyclotomic character is unramified. The restriction of \(f\)
to inertia is a coboundary there, hence is zero.

We need the following form of the rational cyclotomic calculation
\texttt{Kwl\allowbreak{}Onto.\allowbreak{}kwl}. If \(b:G\to\mathbb F_\ell\) is a continuous
additive homomorphism unramified away from an odd prime \(\ell\), then
\(b\) vanishes on the subgroup fixing \(\mu_{\ell^2}\).
Indeed, the Kronecker--Weber theorem with its ramification condition
places the field fixed by \(\ker b\) in \(\mathbb Q(\mu_{\ell^a})\)
for some \(a\ge2\). The unit group modulo \(\ell^a\) is cyclic of
order \(\ell^{a-1}(\ell-1)\). The kernel of reduction modulo
\(\ell^2\) has order \(\ell^{a-2}\). In the cyclic group it consists
of powers of a generator whose exponent is divisible by
\(\ell(\ell-1)\), so each is an \(\ell\)-th power. Every additive
map into \(\mathbb F_\ell\) kills these powers. This proves the claim,
including \(a=2\).

Apply the claim to each composite of \(f\) with a linear functional
\(W[\ell]\to\mathbb F_\ell\). Such functionals separate points, so
\(f\) vanishes on the subgroup fixing \(\mu_{\ell^2}\), hence on
\(\ker\kappa\). The restriction-to-the-cyclotomic-kernel step of
\texttt{Tors\allowbreak{}Class.\allowbreak{}z\allowbreak{}One} now gives \(x=0\). For clarity its injectivity
in this specialization has the following direct verification. Both the
action and the cocycle factor through \(\Gamma=\mathbb Z_\ell^\times\).
For odd \(\ell\), this group has a topological generator \(\gamma\):
its finite factor has order \(\ell-1\), coprime to the orders of all
finite quotients of its procyclic factor \(1+\ell\mathbb Z_\ell\).
Since \(e>0\), the scalar \(\gamma^{-e}-1\) is nonzero.
Choose \(v\in W\) with \((\gamma^{-e}-1)v=f(\gamma)\).
The continuous cocycle \(f-dv\) vanishes on \(\gamma\), hence on all
its integer powers by the cocycle identity, and then everywhere by
density. This proves the vanishing in the first case.

\subsection{The quadratic residual character}

Suppose now that \(\ell-1\mid 2e\) and \(\ell-1\nmid e\).
Write \(n=(\ell-1)/2\). Then \(e=nk\) for an odd integer \(k\).
Since \(e\) is even, \(n\) is even and \(\ell\equiv1\bmod4\).
The residual character \(\overline\kappa^{-e}\) is the quadratic
character of \(\mathbb F_\ell^\times\).

Choose a primitive \(\ell\)-th root \(\zeta\) in \(\Omega\) and let
\(\eta\) be the Legendre symbol, extended by \(\eta(0)=0\).
Set \(s=\sum_{a\bmod\ell}\eta(a)\zeta^a\). Acting on \(\zeta\) by
the exponent \(u\) multiplies \(s\) by \(\eta(u)\), by the change of
variable \(a\mapsto ua\). Replacing \(\zeta\) by \(\zeta^{-1}\)
fixes \(s\) since \(\eta(-1)=1\). Multiplying these two sums and
writing \(a=tb\) in the terms with \(b\ne0\) gives
\[
 s^2
   =\sum_{t\bmod\ell}\eta(t)\sum_{b\ne0}\zeta^{b(t-1)}
   =(\ell-1)-\sum_{t\ne1}\eta(t)=\ell .
\]
Here the sum of the nontrivial character over all residues is zero.
Thus \(K=\mathbb Q(s)=\mathbb Q(\sqrt\ell)\) is quadratic with
discriminant \(\ell\), and its absolute Galois group \(G_K\) is the
kernel of the residual quadratic character. For \(g\in G_K\),
\(\alpha(g)-1\in\ell\mathcal O\). In particular \(G_K\) acts
trivially on \(W[\ell]\).

Choose \(u\in\mathbb Z_\ell^\times\) reducing to a generator of
\(\mathbb F_\ell^\times\). Then \(\overline u^e\ne1\).
At every place \(A\) above \(\ell\) the cyclotomic character maps
inertia onto \(\mathbb Z_\ell^\times\). Indeed the local cyclotomic
extensions of level \(\ell^a\) are totally ramified with full unit
Galois group, as follows by shifting the cyclotomic polynomial by
one to obtain an Eisenstein polynomial. Compactness passes surjectivity
to the inverse limit. Choose \(\delta\) at one such place with
\(\kappa(\delta)=u\). Then \(\alpha(\delta)-1\) is a unit and
\(\delta s=-s\).

Given \(x\in T_{\ell,e}[\ell]\), normalize a representative \(c\)
to have \(c(\delta)=0\), by subtracting
\(d((\alpha(\delta)-1)^{-1}c(\delta))\). Since \(\ell x=0\),
we have \(\ell c=dw\). Evaluation at \(\delta\) gives
\((\alpha(\delta)-1)w=0\), hence \(w=0\). Thus \(c(g)\in W[\ell]\)
for every \(g\). The restriction \(c|_{G_K}\) is an additive
homomorphism. Away from \(\ell\) it vanishes on inertia, since the
character is unramified there and the strict class is locally zero.

At a place \(A\mid\ell\) the strict condition gives \(m_A\in W\)
such that \(c(g)=(\alpha(g)-1)m_A\) for all \(g\) in its
decomposition group. Take \(\tau\) in inertia at this same place
with \(\kappa(\tau)=u\). Then
\((\alpha(\tau)-1)\ell m_A=\ell c(\tau)=0\). The factor is a unit,
so \(m_A\in W[\ell]\). Therefore \(c(g)=0\) for every
\(g\) in the decomposition group that also belongs to \(G_K\).
The restriction to \(G_K\) is unramified at all finite places.
Its target has odd exponent, so it also vanishes on all real inertia.
By the unramified class-field correspondence it factors through the
ideal class group of \(K\).

We give the class-number inequality used here in \texttt{HQuad}.
Put \(X=\sqrt{\ell}\). The quadratic Minkowski bound supplies every
class with an integral ideal of positive norm at most \(X\).
For each integer \(a\ge1\), the number of ideals of norm \(a\) is
at most the number \(d(a)\) of ordered factorizations \(a=ij\).
On a prime power \(p^v\) the split case admits \(v+1\) possible
pairs of exponents above \(p\), and the inert and ramified cases
each admit at most one possibility. Unique factorization of ideals
multiplies these bounds over rational primes. Thus
\[
    h(K)\le\sum_{a\le X}d(a)
      =\sum_{i\le X}\lfloor X/i\rfloor
      \le X(1+\log X)<X^2=\ell .
\]
The harmonic-sum bound follows by integration of \(1/t\), and the
strict inequality from \(\log X<X-1\) for \(X>1\).
Hence \(h(K)\) is coprime to \(\ell\). Every value of \(c|_{G_K}\)
is annihilated both by \(h(K)\) and by \(\ell\), so this restriction
is zero. For \(g\notin G_K\), write \(g=\delta h\), \(h\in G_K\).
Then \(c(g)=c(\delta)+\alpha(\delta)c(h)=0\).
This proves the second case of \eqref{eq:rh-z-vanish}.

The normalization and local restriction just used are the zero-length
specialization of \texttt{Primitive\allowbreak{}Frobenius\allowbreak{}Fixed.\allowbreak{}Sha\allowbreak{}To\allowbreak{}Sh\allowbreak{}T} and
\texttt{Level\allowbreak{}Cohomology.\allowbreak{}Zq\allowbreak{}Field.\allowbreak{}z\allowbreak{}Quad}. In the general-length
formulation the source takes the union over annihilators \(a\ne0\) of
normalized cocycles. Every such annihilator divides a power of \(\ell\)
in this discrete valuation ring. The preceding torsion argument gives
the same conclusion at length zero. The class-field inputs in the
two cases are respectively ramified Kronecker--Weber over \(\mathbb Q\)
and the unramified reciprocity map over the real quadratic field.

\subsection{The estimate at the host level}

In case Z the coefficient prime is odd and outside \(B_\chi\).
Since \(3\in B_\chi\), we have \(\ell\ge5\). The character at level
\(N\) is trivial, and the host also requires
\(\psi(-1)=(-1)^e\). In the characteristic-zero coefficient ring this
forces \(e\) even. The window has \(0<e<t\), so \(e\ge2\).
The last case-Z condition is \(\ell-1\mid2e\); all hypotheses of
\eqref{eq:rh-z-vanish} hold. Passing from the strict condition at
\(\ell\) to the one at \(\ell N\) adds decomposition-group conditions.
It takes a submodule of the group just shown to vanish. The arithmetic
length \(b_{\ell,e}(1)\) of Section~\ref{sec:arithmetic-translation}
is consequently zero at every such host. Its normalized logarithmically
weighted length is zero as well. This proves \texttt{Count\allowbreak{}Z} with the
single exponent \(C_Z=0\), uniformly in all the windows and choices of
coefficients.
